\documentclass[11pt]{article}
\usepackage[letterpaper,margin=1.02in]{geometry}
\usepackage[T1]{fontenc}
\usepackage{lmodern,microtype,amsmath,amssymb,amsthm,mathtools,booktabs,array,enumitem,xcolor}
\usepackage[colorlinks=true,linkcolor=blue!55!black,citecolor=blue!55!black,urlcolor=blue!55!black,pdftitle={The exact quarter-density obstruction for disjoint products of small primes},pdfauthor={Research note for ProveIt}]{hyperref}
\newtheorem{theorem}{Theorem}[section]
\newtheorem{lemma}[theorem]{Lemma}
\newtheorem{corollary}[theorem]{Corollary}
\newtheorem{proposition}[theorem]{Proposition}
\theoremstyle{remark}\newtheorem{remark}[theorem]{Remark}
\newcommand{\Z}{\mathbb Z}\newcommand{\C}{\mathbb C}\newcommand{\Pp}{\mathbb P}\newcommand{\E}{\mathbb E}\newcommand{\floor}[1]{\left\lfloor #1\right\rfloor}\newcommand{\supp}{\operatorname{supp}}\newcommand{\ord}{\operatorname{ord}}\newcommand{\G}{\Gamma}\newcommand{\DGRH}{\textnormal{\scshape DGRH}}
\newcommand{\sourceurl}{https://github.com/VladimirReshetnikov/ProveIt}
\setlist{leftmargin=*,itemsep=2pt,topsep=4pt}
\setlength{\emergencystretch}{2em}
\title{The exact quarter-density obstruction for\\disjoint products of small primes\vspace{3pt}\\\large A finite allocation theorem and its Dirichlet-GRH transfer}
\author{Research note for ProveIt}
\date{27 November 2026}
\begin{document}
\maketitle
\begin{abstract}
The merged ProveIt article on Dirichlet GRH asks for exact support inequalities at the critical retention density $1/4$, where the obstruction quotient for disjoint squarefree products can have order four. We answer that part of the question. For the Klein four-group the answer consists of three linear inequalities coming from quadratic projections; for the cyclic group of order four it consists of four linear inequalities and one conditional parity cut. A closed formula gives the largest number of blocks for which all targets are reachable. In particular, an odd number $r$ of blocks can have every proper marginal fully supported while exactly one projected joint target is missing, although all linear inequalities hold. We prove a general zero-sum packing and completion theorem, giving polynomial-time support decisions in a known bounded quotient at any fixed positive density. The analytic transfer is proved from the Dirichlet GRH statement in the vendored Lean development and a standard integrated explicit formula, with uniform relative counts on occupied fibres. The new finite results are unconditional. We do not independently replay the upstream proof or claim a Lean formalization of the new theorem.
\end{abstract}
\tableofcontents

\section{Input, scope and main result}
\label{sec:input}
The upstream package \cite{lean} is vendored at the pinned upstream commit
\begin{center}\texttt{72fead102f663d5c7a407964a13dc0cf9d977207}.\end{center}
Its \texttt{src/Statements.lean} imports Mathlib and defines the half-plane and critical-strip statements. The relevant exported declarations in \texttt{src/Main.lean} are
\begin{align*}
\texttt{LFunction\_ne\_zero\_of\_half\_lt\_re}&:\ \texttt{DirichletLNonvanishingRightOfHalf},\\
\texttt{LFunction\_eq\_zero\_re}&:\ \texttt{DirichletLZerosOnLine}.
\end{align*}
In mathematical terms, for every Dirichlet character $\chi$ modulo $q\ge1$, $L(s,\chi)$ is nonzero for $\Re s>1/2$, except for the excluded pole of the principal character at $s=1$; every zero in $0<\Re s<1$ lies on $\Re s=1/2$. Primitive and imprimitive characters are included. We call this input \DGRH. It says nothing about simplicity, independence of ordinates, or GRH for general nonabelian or automorphic $L$-functions. No such additional hypothesis is used here.

The proof architecture is relevant: \texttt{New.Reduction.OneSidedBound} obtains a one-sided estimate for a smoothed prime-power remainder from a comparison inequality and a distant cutoff, and exports \texttt{New.grh}; \texttt{New.Main} then uses the functional equation to pass from half-plane nonvanishing to critical-line zeros. Our transfer uses the resulting zero location, not the internal arithmetic bounds. In particular, the smoothing in the downstream character-sum estimate below retains the logarithm-squared prime cutoff, just as the article's small-bias lemma does.

The package README reports a successful build and independent checker replays, using Lean 4.35.0 and Mathlib revision \texttt{56a005855ea17e0b0958f73ec801ae140b0a1679}, with axioms \texttt{propext}, \texttt{Classical.choice}, and \texttt{Quot.sound}. The vendoring README and article explicitly do not recheck the import. We inspected the actual statement and wrapper files, not all hundreds of thousands of proof lines, and did not execute Lean. All number-theoretic implications in this note are therefore stated as implications from \DGRH. The source's reported verification is not our own certification.

\paragraph{Setting.} Let $q\ge3$, $L=\log q$, $G=(\Z/q\Z)^\times$, and
\[
 y=C L^2,\qquad P=\{p\le y:p\text{ prime},\ p\nmid q\},\qquad N=|P|.
\]
An adversary retains a subset $S\subseteq P$, with $m=|S|\ge N/4$. Split $S$ uniformly into labelled blocks $S_0,S_1,\ldots,S_r$ of specified sizes $n_0,n_1,\ldots,n_r$ with $\sum n_j=m$. The reserve is $S_0$; the desired tuple is
\[
 X_j=\prod_{p\in S_j}p\pmod q\quad (1\le j\le r),\qquad h=\min_{0\le j\le r} n_j.
\]
Equivalently, we seek pairwise coprime squarefree positive integers, each with a specified number of prime factors from $S$, in prescribed reduced residue classes. The article \cite{article} (\texttt{dir:thm:disjoint-growing}) gives a bounded obstruction quotient. Its open question \texttt{dir:sec:open}, item ``Density one quarter'', asks for exact inequalities for the order-four cases. We make that classification explicit below and give a separate proof of the transfer in Section~\ref{sec:analytic}.

\begin{theorem}[Quarter-density classification, under \DGRH]\label{thm:main}
Fix $A>0$. There exist constants $C,B>0$ such that for all sufficiently large $q$ and all the above retained sets and quotas with $h\ge BL$, the following hold.
There is a quotient $\pi:G\to\G$ of order $d\le4$, depending only on $S$, not the quotas. Choose a most frequent colour $b\in\G$ of the primes in $S$ and write $\G$ additively. For a target $g\in G^r$ put
\[
 z_j=\pi(g_j)-n_jb,\qquad c_a=\#\{p\in S:\pi(p)-b=a\}\quad(a\in\G).
\]
The target occurs if and only if the corrected vector $(z_j)$ passes the corresponding test in Section~\ref{sec:tests}. In particular:
\begin{enumerate}[label=(\roman*)]
\item If $\G=C_4$, write $(u,v,w)=(c_1,c_2,c_3)$ and $(i,j,k)=(\#\{z=1\},\#\{z=2\},\#\{z=3\})$. The four inequalities
\begin{align}
 u+w&\ge i+k,\label{eq:odd}\\
 3u+2v+w&\ge 3i+2j+k,\label{eq:cost1}\\
 u+2v+3w&\ge i+2j+3k,\label{eq:cost3}\\
 u+2v+w&\ge i+2j+k\label{eq:cost2}
\end{align}
and, only when $j>v$, $i\le u$, $k\le w$, the additional inequality
\begin{equation}\label{eq:parity}
 \floor{(u-i)/2}+\floor{(w-k)/2}\ge j-v
\end{equation}
are necessary and sufficient.
\item If $\G=C_2\times C_2$, label its nonzero elements $1,2,3$ with $1+2=3$. The necessary and sufficient inequalities are
\begin{equation}\label{eq:klein}
 u+v\ge i+j,\qquad u+w\ge i+k,\qquad v+w\ge j+k.
\end{equation}
\end{enumerate}
For every occupied target, if $\nu$ denotes the law of $(\pi(X_1),\ldots,\pi(X_r))$, then
\begin{equation}\label{eq:lift}
 \Pp(X=g)=\frac{d^r}{|G|^r}\,\nu(\pi^r(g))(1+\theta_g),\qquad |\theta_g|\le q^{-A}.
\end{equation}
These statements hold simultaneously for every admissible number of blocks; in particular $r$ can grow on the order of $L/\log L$ with balanced quotas.
\end{theorem}

The two order-four tests differ qualitatively: the Klein case is completely detected by its three quadratic projections, whereas the cyclic case requires an integral constraint even after all its listed real linear inequalities pass.

\section{A general zero-sum completion theorem}\label{sec:packing}
We prove a finite statement independent of number theory. An item is labelled, and its colour is an element of a finite abelian group $\G$ written additively. Translate all colours by a chosen $b$, so the baseline is $0$. A corrected target for a block of size $n_j$ is its original target minus $n_jb$.

\begin{lemma}[Short witnesses]\label{lem:short}
Every multiset of nonzero colours summing to $z$ has a submultiset summing to $z$ of length at most $d-1$, where $d=|\G|$. For $z=0$ the empty multiset suffices.
\end{lemma}
\begin{proof}
Delete nonempty zero-sum submultisets until no such submultiset remains. If a remaining sequence has length at least $d$, its first $d+1$ partial sums (including zero) contain a repeated value, giving a nonempty consecutive zero-sum subsequence. This is a contradiction. The deletion procedure preserves the sum.
\end{proof}

\begin{theorem}[Exact quota-free reduction]\label{thm:packing}
Let $\G$ have order $d$, and let $m$ items have translated colour counts $(c_a)_{a\in\G}$. There are $r$ recorded blocks and one reserve, with quotas $n_0,\ldots,n_r$ summing to $m$. Assume
\begin{align}
 n_j&\ge d-1\quad(1\le j\le r),\label{eq:quota1}\\
 n_0&\ge M_\G:=\sum_{a\ne0}(\ord(a)-1),\label{eq:quota2}\\
 c_0&\ge (d-1)(r+1).\label{eq:quota3}
\end{align}
Then the following are equivalent:
\begin{enumerate}[label=(\alph*)]
\item There is a partition of all items into these quotas whose recorded sums are the corrected targets $z_1,\ldots,z_r$.
\item There exist disjoint multisets of nonzero-colour items, one for each target, each summing to that target, with no quotas imposed and unused items allowed.
\item The multisets in (b) can be required to have length at most $d-1$.
\end{enumerate}
The implication (c) to (a) is constructive; no search over complete partitions is necessary.
\end{theorem}
\begin{proof}
Deleting zero-colour items gives (a) to (b); Lemma~\ref{lem:short} gives (b) to (c). For (c) to (a), place the witnesses in their recorded blocks; \eqref{eq:quota1} ensures they fit. For each nonzero colour $a$, split its remaining items into packets of size $\ord(a)$ and one leftover group of size less than $\ord(a)$. Each full packet has sum zero. Put all leftovers in the reserve, possible by \eqref{eq:quota2}.

Now distribute the full packets, in any order, to any block with enough remaining capacity. If a packet of size $s\le d$ fails to fit, the sum of remaining capacities is at most $(s-1)(r+1)\le(d-1)(r+1)$. But that sum is exactly the number of unplaced items, including all $c_0$ zero-colour items and at least the pending packet of size $s>0$. This contradicts \eqref{eq:quota3}. Thus all packets fit. Fill the remaining slots with zero-colour items. No recorded target sum has changed. This also works for $d=1$, when there are no nonzero colours.
\end{proof}

\begin{corollary}[Polynomial-time test in a known quotient]\label{cor:dp}
For fixed $d$, under \eqref{eq:quota1}--\eqref{eq:quota3}, feasibility can be decided in $O_d\bigl(r(1+(d-1)r)^{d-1}\bigr)$ arithmetic operations on small integers, once the colour counts and targets are supplied. A witness can be reconstructed by storing predecessors and then applying Theorem~\ref{thm:packing}.
\end{corollary}
\begin{proof}
For each $z\in\G$, list all vectors $e\in\Z_{\ge0}^{d-1}$ with $\sum_a e_a\le d-1$ and $\sum_a e_a a=z$; this is a constant-size list when $d$ is fixed. For $z=0$ use only the zero vector. Cap each available nonzero count at $(d-1)r$, since at most that many items can be used in short witnesses. Start with the all-zero used-count vector. For each recorded target, add all its candidate vectors to previously reachable vectors and discard sums exceeding the capped supply. The number of possible states is at most $(1+(d-1)r)^{d-1}$. The final state set is nonempty exactly when short witnesses exist.
\end{proof}

The same construction applies to every fixed retention fraction $a>0$ in place of $1/4$, since Section~\ref{sec:analytic} gives $d\le\floor{1/a}$. Indeed $c_0\ge m/d$, $n_0,h\to\infty$, and $m\ge(r+1)h$, so the three finite hypotheses hold for sufficiently large $q$. This is a polynomial-time test \emph{given the quotient}; it does not solve the article's separate problem of finding that quotient in time polynomial in $\log q$.

\begin{lemma}[Direct matches can always be prioritized]\label{lem:direct}
Let $s_a$ be the supply of nonzero colour $a$ and $t_a$ the number of nonzero targets of that colour. In the quota-free problem (b) of Theorem~\ref{thm:packing}, if a feasible collection exists, one exists with $\min(s_a,t_a)$ singleton witnesses of colour $a$ for every $a\ne0$. Therefore only deficient target types need to be formed from surplus item types after direct matching.
\end{lemma}
\begin{proof}
Among all feasible collections choose one maximizing the total number of singleton witnesses. Take all zero targets empty. Suppose the number of singleton matches of colour $a$ is less than $\min(s_a,t_a)$. Some non-singleton witness $W$ has sum $a$, and an item $x$ of colour $a$ is not used in any singleton witness. If $x$ is unused, replace $W$ by $\{x\}$. If it is in $W$, retain only $x$. If it is in a different non-singleton witness $V$, replace $W$ by $\{x\}$ and replace $V$ by $(V\setminus\{x\})\cup W$. The latter witness still has its original sum, since $W$ sums to $a$. Disjointness holds in every case. No old singleton is lost and a new one is gained, contradicting maximality. Thus all direct matches can be made. Delete any zero-sum submultisets from remaining witnesses if necessary. Only surplus colours remain for these witnesses.
\end{proof}

\section{Complete support criteria in groups of order at most four}\label{sec:tests}
Throughout this section, the finite quota hypotheses of Theorem~\ref{thm:packing} are in force. We prove the quota-free criteria and then transfer them back. For all cyclic groups, coordinates are labelled by their natural residues.

\subsection{Groups of order one, two and three}
For the trivial group every target is reachable. For $C_2$, if $u$ is the number of nonzero items and $i$ the number of nonzero targets, the exact criterion is $u\ge i$. For $C_3$, let supplies be $u,v$ and demands $i,j$ of colours $1,2$. The exact criterion is
\begin{equation}\label{eq:c3}
 2u+v\ge2i+j,\qquad u+2v\ge i+2j.
\end{equation}
For necessity, a minimal nonzero witness of sum $1$ is either a single $1$ or two $2$'s, and analogously for sum $2$. Apply the weights $(2,1)$ and $(1,2)$. For sufficiency, use direct matching. Both types cannot be deficient by either inequality. If $u<i$, each remaining target $1$ needs two $2$'s, and $v-j\ge2(i-u)$ is precisely the first inequality. The case $v<j$ is symmetric. Otherwise direct matching suffices. These recover the previously established order-two and order-three tests.

\subsection{The Klein four-group}
Label its nonzero elements $1,2,3$ with $1+2=3$. A minimal witness for each nonzero sum is a singleton of that colour or the two other colours. Indeed equal colours cancel in pairs, and all three distinct nonzero elements sum to zero.

\begin{theorem}[Klein criterion, unconditional]\label{thm:klein}
For supplies $(u,v,w)$ and demands $(i,j,k)$ the inequalities \eqref{eq:klein} are necessary and sufficient.
\end{theorem}
\begin{proof}
Each pair sum of supplies counts the nonzero items seen by one of the three nontrivial homomorphisms to $C_2$. Each demanded target in that nonzero fibre needs at least one such item, giving necessity. For sufficiency, two deficient colours are impossible by their pair inequality. If no colour is deficient, direct matching is enough. If, for example, $u<i$, the first and second inequalities give $v-j\ge i-u$ and $w-k\ge i-u$. After matching directly, pair one surplus $2$ and one surplus $3$ for each deficient $1$. The other possibilities are symmetric.
\end{proof}

\subsection{The cyclic four-group}
The following table lists every nonempty zero-sum-free multiset with specified nonzero sum in $C_4$. It follows either by Lemma~\ref{lem:short} and enumerating lengths $1,2,3$, or by observing that $1+3$ and $2+2$ are forbidden zero-sum pairs.
\begin{center}
\begin{tabular}{cl}
\toprule
Target & Minimal witnesses \\
\midrule
$1$ & $\{1\},\ \{2,3\},\ \{3,3,3\}$\\
$2$ & $\{2\},\ \{1,1\},\ \{3,3\}$\\
$3$ & $\{3\},\ \{1,2\},\ \{1,1,1\}$\\
\bottomrule
\end{tabular}
\end{center}

\begin{theorem}[Cyclic criterion, unconditional]\label{thm:cyclic}
The four linear inequalities \eqref{eq:odd}--\eqref{eq:cost2}, together with the conditional floor inequality \eqref{eq:parity}, are necessary and sufficient for reachability in $C_4$.
\end{theorem}
\begin{proof}
\emph{Necessity.} Assign costs to colours $1,2,3$ using any of the four vectors
\[
 (1,0,1),\quad(3,2,1),\quad(1,2,3),\quad(1,2,1).
\]
Inspection of the witness table shows that every witness of a given sum costs at least the cost of that sum. All costs are nonnegative. Summing costs over disjoint witnesses gives the four linear inequalities. If $j>v$, $i\le u$, $k\le w$, Lemma~\ref{lem:direct} permits direct matches for all odd targets and all $v$ available items of colour $2$. The remaining $j-v$ targets of colour $2$ must each use a pair of equal odd colours. There are at most $\floor{(u-i)/2}+\floor{(w-k)/2}$ such disjoint pairs. This proves the additional condition.

\emph{Sufficiency.} First match directly as many targets as possible. Write $U=u-i$, $V=v-j$, $W=w-k$ for signed surpluses. We construct witnesses for every case of negative coordinates.
\begin{enumerate}[label=\textnormal{Case \arabic*.}]
\item No negative coordinate: direct matches suffice.
\item Only $U<0$. Put $D=-U>0$. By \eqref{eq:odd} and \eqref{eq:cost1}, $W\ge D$ and $2V+W\ge3D$. Make $x=\min(V,D)$ remaining targets of sum $1$ from a $2$ and a $3$; make the remaining $D-x$ from triples of $3$'s. The consumption of colour $3$ is $3D-2x\le W$, using the two displayed bounds according as $V\ge D$ or $V<D$. The case where only $W<0$ is symmetric.
\item Only $V<0$. Inequality \eqref{eq:parity} gives enough equal-odd-colour pairs. Use any $-V$ such pairs from the two available types.
\item Two negative coordinates. If $U,W<0$, \eqref{eq:odd} fails. If $V,W<0$, then \eqref{eq:cost3} gives $U\ge 2(-V)+3(-W)$: form the deficient $2$'s from pairs of $1$'s and the deficient $3$'s from triples of $1$'s. If $U,V<0$, reverse $1$ and $3$ and use \eqref{eq:cost1}. Three negative coordinates are ruled out by \eqref{eq:odd}.
\end{enumerate}
All allocations are disjoint by their supply bounds. Theorem~\ref{thm:packing} completes them to the prescribed quotas.
\end{proof}

The floor inequality is a genuinely integral cut. In its active case put $D=j-v$, $U=u-i\ge0$, $W=w-k\ge0$. The fourth linear inequality is $U+W\ge2D$. Given this bound, the floor inequality fails exactly when $U,W$ are both odd and $U+W=2D$. Thus all real-linear tests can pass while the target is unreachable. Already $(u,v,w)=(1,0,1)$ and $(i,j,k)=(0,1,0)$ demonstrate this phenomenon.

\section{Full coverage, exact support size and sharp examples}\label{sec:coverage}
Define the coverage capacity of a colour multiset as the largest number of recorded targets for which every corrected vector is reachable in the quota-free problem. If the group is trivial the capacity is infinite. In all statements on actual partitions, quotas must meet Theorem~\ref{thm:packing}.

\begin{theorem}[Exact coverage capacities, unconditional]\label{thm:capacity}
The capacities for the nontrivial groups of order at most four are
\begin{center}
\begin{tabular}{@{}ll@{}}
\toprule
Group & Capacity $R_\G$ (supply coordinates as above)\\
\midrule
$C_2$ & $u$\\
$C_3$ & $\floor{\min(2u+v,u+2v)/2}$\\
$C_2\times C_2$ & $\min(u+v,u+w,v+w)$\\
$C_4$ & $\min\left\{u+w,\floor{\frac{3u+2v+w}{3}},\floor{\frac{u+2v+3w}{3}},\ v+\floor{u/2}+\floor{w/2}\right\}$\\
\bottomrule
\end{tabular}
\end{center}
All $r$-tuples occur exactly when $r$ does not exceed the listed capacity. Testing the constant corrected target vectors alone suffices to decide this property.
\end{theorem}
\begin{proof}
For $C_2,C_3,C_2\times C_2$, maximize the right sides of the linear tests over nonnegative target counts with total at most $r$. A maximizer is always a constant nonzero target vector, and the resulting bounds are the listed ones.

For $C_4$, the constant vectors of sums $1$ and $3$ force $u+w\ge r$, $3u+2v+w\ge3r$ and $u+2v+3w\ge3r$. The constant vector of sum $2$ forces $v+\floor{u/2}+\floor{w/2}\ge r$ by the direct-match lemma and witness table. Conversely, assume these four bounds. The first three linear tests hold for all target histograms because $i+j+k\le r$. The fourth holds since $u+2v+w\ge2r\ge i+2j+k$. If the conditional floor test is active, then
\begin{align*}
 \floor{(u-i)/2}+\floor{(w-k)/2}
 &\ge \floor{u/2}+\floor{w/2}-i-k\\
 &\ge r-v-i-k\ge j-v.
\end{align*}
Theorem~\ref{thm:cyclic} completes the argument.
\end{proof}

\begin{corollary}[Exact cardinality]
In either group of order four let $\mathcal H_r$ consist of the triples $(i,j,k)$ of nonnegative integers with $i+j+k\le r$ that pass its support test. Under the analytic transfer, the exact number of reachable residue tuples is
\begin{equation}\label{eq:card}
 \left(\frac{|G|}{4}\right)^r
 \sum_{(i,j,k)\in\mathcal H_r}\frac{r!}{i!j!k!(r-i-j-k)!}.
\end{equation}
This sum has at most $\binom{r+3}{3}$ terms. Analogous formulas hold for smaller quotients.
\end{corollary}
\begin{proof}
Each histogram corresponds to the indicated number of ordered corrected quotient targets; each corrected target is a translate of one ordinary quotient target. Every occupied target has exactly $(|G|/4)^r$ lifts, all of which are attainable by \eqref{eq:lift}. The formula follows.
\end{proof}

\begin{proposition}[A single missing pattern at arbitrarily high arity]\label{prop:missing}
Let $r$ be odd and take cyclic supplies $u=w=r$, $v=0$. Exactly one corrected quotient vector is missing: the vector with every entry equal to $2$. In particular every proper marginal has full support, and all four linear inequalities hold even for the missing vector.
\end{proposition}
\begin{proof}
The four linear tests hold for every histogram of total at most $r$. The only potential failure is the floor cut when $j>0$. Since $r$ is odd,
\[
 \floor{(r-i)/2}+\floor{(r-k)/2}=r-1-\floor{i/2}-\floor{k/2}.
\]
If $i+k\ge1$, then $i+k-\floor{i/2}-\floor{k/2}\ge1$, so this is at least $r-i-k\ge j$. If $i=k=0$, the floor cut holds if and only if $j\le r-1$. Thus precisely $(i,j,k)=(0,r,0)$ fails. A proper marginal can always be extended by setting an unrecorded coordinate to zero.
\end{proof}

\paragraph{Testing without selecting an order-four quotient.}
At density $1/4$, it is sufficient (and necessary) to test every Dirichlet character of exact order $2,3,4$ separately, using its cyclic quotient, its own most frequent colour, and the corresponding criteria above. Necessity follows by projection, with the finite quota conditions automatic. For sufficiency take the actual obstruction quotient $\G$. If it is cyclic, the generating character is among the tests. If it is Klein, the three pair inequalities are precisely its three quadratic projections with baseline the image of the chosen modal $b\in\G$. If such an image is not a modal colour in a particular quadratic projection, at least $m/2\ge r$ items are exceptional relative to it, so its inequality holds automatically. If it is modal, the corresponding quadratic test applies (ties are also harmless). Therefore all Klein inequalities hold. This characterization uses only low-order characters, but their enumeration for arbitrary composite $q$ is not claimed to be polynomial in $\log q$.

\section{The analytic transfer from Dirichlet GRH}\label{sec:analytic}
We rederive the bounded-quotient lifting argument used in the article \cite{article}, including the lower bound for occupied quotient fibres, so the new conclusion does not rely on an unchecked allocation theorem in that article. The only deep analytic input is \DGRH\ and the standard explicit formula for Dirichlet $L$-functions. No new assertion about zero statistics is needed.

\begin{lemma}[Small bias at the logarithm-squared cutoff, under \DGRH]\label{lem:bias}
For any fixed $\epsilon>0$ there is $C>0$ such that for every sufficiently large $q$, the above full pool $P$ satisfies
\[
 \max_{\chi\ne1}\left|\frac1N\sum_{p\in P}\chi(p)\right|\le\epsilon,
 \qquad N\sim\frac{CL^2}{\log(CL^2)}.
\]
\end{lemma}
\begin{proof}
The once-integrated GRH explicit-formula estimate for every nonprincipal character modulo $q$ is
\begin{equation}\label{eq:integrated}
 \left|\sum_{n\le t}(t-n)\Lambda(n)\chi(n)\right|\ll t^{3/2}\log(2q),\qquad t\ge2,
\end{equation}
with an absolute implied constant; see Montgomery--Vaughan \cite[Chapter 13, (13.28)]{mv} and the article's \texttt{dir:lem:bias}. The kernel $1/(s(s+1))$ makes the zero sum absolutely convergent, and the unit-interval zero-count bound $O(\log(q(2+|T|)))$ gives the factor $\log(2q)$. For an induced character, omitted prime-power terms have total integrated weight at most $t\sum_{p\mid q}\log p\,\lfloor\log t/\log p\rfloor\ll t\log t\log(2q)$, absorbed by the right side. This also shows why no uniform sharp-cutoff GRH estimate with a missing logarithm is needed.

Fix $0<\eta<1/2$ and put
\[
 W(n)=\frac{(y-n)_+-((1-\eta)y-n)_+}{\eta y}.
\]
It equals one up to $(1-\eta)y$ and decreases linearly to zero. Subtract two instances of \eqref{eq:integrated} and divide by $\eta y$ to bound the weighted von Mangoldt sum by $O(\sqrt y\log(2q)/\eta)$. Removing prime powers costs $O(\sqrt y\log^2 y)$; replacing $W$ by the sharp cutoff costs at most $\vartheta(y)-\vartheta((1-\eta)y)=(\eta+o(1))y$. Finally replacing each prime weight $\log p$ by $\log y$ costs at most $\pi(y)\log y-\vartheta(y)=o(y)$. These last assertions use the ordinary prime number theorem. Primes dividing $q$ number at most $L/\log2$, so $N\sim y/\log y$ and all errors are uniform in the character. Divide by $N\log y\sim y$ to obtain a bias at most
\[
 \eta+O\left(\frac1{\eta\sqrt C}\right)+o(1).
\]
Choose $\eta$ sufficiently small and then $C$ sufficiently large. This proves the assertion.
\end{proof}

\begin{lemma}[The obstruction quotient]\label{lem:quotient}
Fix $0<a\le1$ and $D=\floor{1/a}$. For sufficiently small $\tau,\epsilon>0$ depending only on $a$, if $P$ has bias at most $\epsilon$ as in Lemma~\ref{lem:bias} and $|S|/N=\alpha\ge a$, the subgroup $H\le\widehat G$ generated by characters with
\[
 b_\chi=\left|m^{-1}\sum_{p\in S}\chi(p)\right|>1-\tau
\]
has order at most $D$. All characters outside $H$ have bias at most $1-\tau$ on $S$. The natural evaluation quotient $\pi:G\to\widehat H=\G$ is onto.
\end{lemma}
\begin{proof}
The Gram matrix of $t$ distinct characters restricted to $P$, for the normalized counting inner product, has diagonal one and off-diagonal entries of modulus at most $\epsilon$. Its largest eigenvalue is at most $1+(t-1)\epsilon$. Taking inner products with the indicator $1_S$ of squared norm $\alpha$ yields
\begin{equation}\label{eq:gram}
 \alpha^2\sum_{\chi\ \text{among the }t}b_\chi^2\le \alpha(1+(t-1)\epsilon).
\end{equation}
A generator of $H$ is within $\sqrt{2\tau}$ in normalized $L^2(S)$ of some unimodular constant, and so is its inverse. By telescoping products of unimodular functions, a word of length at most $D$ is within $D\sqrt{2\tau}$ of a constant and has bias at least $1-D^2\tau$. If $H$ contains at least $D+1$ elements, the word ball of radius $D$ contains at least $D+1$ (a connected Cayley graph expands until it is exhausted). Taking these characters in \eqref{eq:gram} contradicts
\[
 (D+1)a(1-D^2\tau)^2>1+D\epsilon.
\]
Such constants exist since $(D+1)a>1$; for $a=1/4$, $\tau=1/256$, $\epsilon=1/100$ work. The outside bias bound is by definition, and the evaluation map is onto by duality of finite abelian groups (its kernel is the annihilator of $H$).
\end{proof}

\begin{lemma}[Fourier decay across quotient cosets]\label{lem:fourier}
Let $k_0=r+1$, $\alpha_j=n_j/m$, and fix a character tuple $(\chi_1,\ldots,\chi_r)$ with $\chi_0=1$. If $t$ of its $k_0$ labels lie outside a coset of $H$ of largest total quota weight, then
\begin{equation}\label{eq:fourier}
 \left|\E\prod_{j=1}^r\chi_j(X_j)\right|\le(m+1)^{k_0}e^{-\tau th/2}.
\end{equation}
\end{lemma}
\begin{proof}
Assign independently to each prime a label $j$ with probability $\alpha_j$, then condition on the quota counts. The probability $p_n$ of the specified counts is at least $(m+1)^{-k_0}$: the vector $n$ is a multinomial mode, since transferring one count from an overfull label $i$ to an underfull label $j$ increases the probability by the factor $v_i\alpha_j/((v_j+1)\alpha_i)\ge1$. There are at most $(m+1)^{k_0}$ possible count vectors.

The unnormalized conditional expectation is the coefficient of $Z_0^{n_0}\cdots Z_r^{n_r}$ in
\[
 \prod_{p\in S}\left(\sum_{j=0}^r\alpha_j Z_j\chi_j(p)\right).
\]
Bound the coefficient by the maximum modulus on the unit torus. For any unimodular $v_j$,
\[
 \left|\sum_j\alpha_j v_j\right|^2
 =1-2\sum_{i<j}\alpha_i\alpha_j(1-\Re(v_i\overline v_j)).
\]
Use $1-x\le e^{-x}$, take square roots, and multiply over all primes. For labels in distinct $H$-cosets, the mean of $\chi_i\overline{\chi_j}$ on $S$ has modulus at most $1-\tau$. Thus the modulus of the product is at most
\[
 \exp\left(-m\tau\sum_{\text{different cosets }i<j}\alpha_i\alpha_j\right).
\]
If the coset masses are $w_C$ and their maximum is $M$, the sum is $(1-\sum_C w_C^2)/2\ge(1-M)/2\ge th/(2m)$. Divide by $p_n$ to obtain \eqref{eq:fourier}.
\end{proof}

\begin{lemma}[Occupied projections have adequate mass]\label{lem:repair}
For $m$ labelled items of colours in a finite abelian group of order $d$ and exponent $e$, allocated uniformly into $k_0$ blocks of fixed quotas, the probability of every attainable vector of recorded block sums is at least $(1+mk_0)^{-6edk_0}$.
\end{lemma}
\begin{proof}
Choose a reference allocation attaining the vector, including its reserve sum. Write its table of colour counts as $Y$, and let $X$ be the table of any allocation. The tables have identical row and column margins. Let $\sigma_{ij}$ be the least nonnegative residue of $Y_{ij}$ modulo $e$. Set
\[
 Z_{ij}=\max\{0,\floor{(X_{ij}-\sigma_{ij})/e}\}.
\]
The table $\sigma+eZ$ differs from $X$ by less than $e$ in each entry. The desired margins for $Z$ are those of $(Y-\sigma)/e$; each family of margin errors has total absolute size at most $dk_0$. Remove units from overfull rows and then overfull columns, at a cost at most $dk_0$ each. All remaining margins are deficient or correct, with total common deficit at most $3dk_0$. Add units at intersections of deficient rows and columns until margins match, which is possible since both deficits have the same total. The resulting nonnegative table $Z'$ differs from $Z$ in entrywise $\ell^1$ norm by at most $5dk_0$.

The table $X'=\sigma+eZ'$ therefore has the original margins, is congruent to $Y$ entrywise modulo $e$, and differs from $X$ in $\ell^1$ norm by at most $6edk_0$. The congruence ensures the required colour sums because $ea=0$ for every colour. Move at most $6edk_0$ individual item labels to realize $X'$. Every allocation is within that Hamming distance of an attaining allocation; a Hamming ball has at most $(1+mk_0)^{6edk_0}$ allocations, by coding each changed item and destination with padding. A covering argument gives the stated probability.
\end{proof}

\begin{theorem}[Uniform lifting at arbitrary fixed retention density, under \DGRH]\label{thm:lift}
Fix $a>0$, $a\le1$, and $A>0$. There are $C,B>0$ such that the support identity and relative count estimate \eqref{eq:lift} hold whenever $|S|\ge aN$ and $h\ge BL$, for all large $q$. The quotient has order $d\le D=\floor{1/a}$. Moreover, the exact support is decided in the quotient by Theorem~\ref{thm:packing} and Corollary~\ref{cor:dp}.
\end{theorem}
\begin{proof}
Choose $\tau,\epsilon,C$ by Lemmas~\ref{lem:bias} and \ref{lem:quotient}. The characters in $H^r$ contribute exactly $d^r\nu(\pi^r(g))$ to the Fourier inversion expression for $|G|^r\Pp(X=g)$. For a tuple outside $H^r$, the number $t$ of labels outside the largest coset is at least one, because the reserve has character $1$. For given $t$, count at most $q\binom{k_0}t d^{k_0}q^t$ tuples (choose a coset, the outside labels, then their characters and the inside characters). By Lemma~\ref{lem:fourier} the total absolute error is at most
\begin{equation}\label{eq:error}
 qd^{k_0}(m+1)^{k_0}\left[(1+qe^{-\tau h/2})^{k_0}-1\right].
\end{equation}
For $B\ge1$ and large $q$, $k_0 h\le m\le N\sim CL^2/\log(CL^2)$, hence $k_0\log(m+1)\le K L$ for $K=3(C+1)$; also $k_0\le q$, $d\le m+1$. Take $\tau B/2>3$ so the bracket is at most $2k_0qe^{-\tau h/2}$. Thus \eqref{eq:error} is at most $2q^{2K+3-\tau B/2}$.

By Lemma~\ref{lem:repair}, any positive $\nu$ is at least $(1+mk_0)^{-6D^2k_0}\ge q^{-12D^2K}$. Taking $B$ large enough that $\tau B/2\ge 2K+12D^2K+A+5$ bounds the absolute error by $q^{-A}$ times $d^r\nu$ on all occupied points. The probability is then positive, and an unoccupied projection has no lift by definition. This proves \eqref{eq:lift} and the support identity. Finally the most frequent colour satisfies $c_0\ge m/d\ge k_0h/d$, and all quotas are at least $h\to\infty$, so \eqref{eq:quota1}--\eqref{eq:quota3} hold. The finite theorem applies.
\end{proof}

Taking $a=1/4$ and applying Section~\ref{sec:tests} proves Theorem~\ref{thm:main}. Balanced quotas exist with $k_0$ of order $L/\log L$, since $m\ge aN$. No larger order of $r$ can meet the same logarithmic quota requirement with this pool size. The constants here are existence constants, not numerically optimized; the result holds for all sufficiently large moduli.

\section{Arithmetic sharpness and a sparse-exception law}\label{sec:examples}
\begin{corollary}[Actual arithmetic example of the parity cut, under \DGRH]\label{cor:primeexample}
For all sufficiently large primes $q\equiv1\pmod4$, there is a retained set of exactly $\lceil N/4\rceil$ primes and admissible quotas with any odd $r\ge1$ in a range $r\le cL/\log L$ (for a fixed sufficiently small $c>0$), such that all proper marginals of $X$ are fully supported and precisely $(|G|/4)^r$ tuples are missing. The omitted tuples form one Cartesian product of quartic residue classes, and all four linear inequalities pass for their projected targets.
\end{corollary}
\begin{proof}
Choose an order-four character of $G$, available because $G$ is cyclic of order divisible by four. Its four colour classes in $P$ each have size at least $(1-3\epsilon)N/4$ by Fourier inversion and Lemma~\ref{lem:bias}. A largest class has at least $\lceil N/4\rceil$ items. Retain $\lceil N/4\rceil-2r$ items of that colour $b$, and $r$ items of each adjacent quartic colour $b+1,b+3$; retain none of colour $b+2$. For large $q$ all these selections are possible, $2r/m\to0$, and $b$ is the unique modal colour. The quartic character has bias at least $1-4r/m>1-\tau$ on $S$. Hence $H$ contains a cyclic subgroup of order four and, by Lemma~\ref{lem:quotient}, is exactly that subgroup. Choose balanced quotas with $h\ge BL$, possible for the stated range of $r$. Proposition~\ref{prop:missing}, together with Theorem~\ref{thm:lift}, gives all assertions. In a proper marginal any given residue tuple can be extended by taking one omitted coordinate with corrected quotient colour zero; such a full target occurs.
\end{proof}

This example shows why merely extending the real linear inequalities from the order-three case cannot settle the order-four question. It is a true arithmetic obstruction at exactly the threshold, not just a finite toy model. Primes $q\equiv1\pmod4$ are unbounded by Dirichlet's theorem.

\begin{proposition}[Sparse exceptions, uniform in the number of blocks]\label{prop:iid}
Let $s=m-c_0$ be the number of exceptional items in a quotient of any fixed order $d$. Let $\widetilde\nu$ be the corrected quotient law obtained by putting each exceptional item independently into a label $j$ with probability $n_j/m$, ignoring the baseline zero-colour items. Then the corrected actual quotient law $\nu_{\rm corr}$ satisfies
\begin{equation}\label{eq:tv}
 d_{\rm TV}(\nu_{\rm corr},\widetilde\nu)\le 1-\frac{(m)_s}{m^s}\le \frac{s(s-1)}{2m},
\end{equation}
where $(m)_s=m(m-1)\cdots(m-s+1)$. In particular, the bound is uniform in $r$, including growing $r$.
\end{proposition}
\begin{proof}
Regard each label as a collection of $n_j$ slots. The positions of the $s$ labelled exceptions in a uniform allocation have the law of $s$ independent uniform slots conditioned on being distinct. The total variation distance of an unconditional distribution from its conditioning on an event of positive probability equals the probability of the complement. The probability of at least one collision is at most $\binom{s}{2}/m$. Map slot positions to their block sums; total variation cannot increase. This proves both bounds.
\end{proof}

In the independent law, the Fourier coefficient at a character tuple $\boldsymbol\psi\in\widehat\G^r$ is exactly
\[
 \prod_{a\ne0}\left(\frac{n_0}{m}+\sum_{j=1}^r\frac{n_j}{m}\psi_j(a)\right)^{c_a}.
\]
Thus the counts $c_a$ and relative quotas completely specify the leading probability profile in the regime $s=o(\sqrt m)$, without keeping individual prime identities. The relative error of Theorem~\ref{thm:lift} also implies total variation distance at most $q^{-A}/2$ between the full law and uniform lifting of its actual quotient law. Consequently uniform lifting of the independent-exception model differs in total variation by at most $q^{-A}/2+s(s-1)/(2m)$ after undoing the translation. For Corollary~\ref{cor:primeexample}, $s=2r=O(L/\log L)$ and $m\asymp L^2/\log L$, so the latter error is $O(1/\log L)$. This addresses sparse-exception probability profiles but not the article's entire open problem concerning all growing-dimensional mixing regimes.

\section{Verification and novelty status}
The supplied \texttt{verify.py} uses integer arithmetic and no third-party dependencies. Its principal oracle processes one available item at a time, sending it to a recorded block or the reserve and tracking histograms of current sums. It does not use the inequalities or the minimal-witness table. A second oracle enumerates colour-count tables of specified quotas, independently of the first oracle and the general bounded-quotient dynamic program. The logged run checked:
\begin{itemize}
\item 479,682 exact support comparisons for $C_4$ and $C_2\times C_2$, each supply coordinate from 0 through 8, and 1 through 7 recorded blocks;
\item 10,206 exact full-coverage comparisons over the same range;
\item 5,280 fixed-quota target comparisons, also checked against the general dynamic program, and the single-missing-pattern family for odd $r$ from 1 to 23.
\end{itemize}
These computations are regression tests, not substitutes for the proofs. The finite proof can be formalized using finite multisets, partial sums, and elementary natural-number inequalities; an outline and upstream interface audit file are included, explicitly uncompiled.

\paragraph{What is new relative to the inspected article.}
The article's \texttt{dir:thm:disjoint-growing} already supplies the analytic quotient law, and \texttt{dir:thm:disjoint-support} handles orders two and three. We claim neither as new. The new contribution is the quota-free completion lemma, direct-match normal form, exact order-four inequalities with their parity cut, capacity and support-size formulas, general bounded-quotient decision procedure, and the arithmetic single-missing-pattern application. The quarter-density question in \texttt{dir:sec:open} expressly asks for these order-four support inequalities. A focused literature search found general work on products of small primes in arithmetic progressions \cite{products} and on zero-sum-free sequences \cite{zerosum}, but did not establish a prior identical classification. The result is proposed as a new contribution to the repository, not a certified priority claim across all literature. No assertion is made that the entire 770-page article has been verified, or that its other open questions are resolved.

\begin{thebibliography}{9}
\bibitem{lean} L. Sernau et al., \emph{The Generalized Riemann Hypothesis for Dirichlet L-functions, in Lean 4}, readable edition, upstream commit 72fead102f663d5c7a407964a13dc0cf9d977207, 25 November 2026 (as identified in the vendor README). \href{https://github.com/lsernau/grh-lean/tree/72fead102f663d5c7a407964a13dc0cf9d977207/grh}{Upstream}. \href{https://github.com/VladimirReshetnikov/ProveIt/tree/main/Analysis/GeneralizedRiemannHypothesis/vendor/grh-lean}{Vendored source}.
\bibitem{article} ProveIt Contributors, \emph{Beyond the Riemann hypothesis for Dirichlet L-functions: consequences, transfers and limits}, 26 November 2026, particularly \texttt{sections/src.tex}, \texttt{sections/dir.tex}. \href{https://github.com/VladimirReshetnikov/ProveIt/tree/main/Analysis/GeneralizedRiemannHypothesis/Article}{Article source}. Inspected \texttt{dir.tex} blob: \texttt{5b090c0707c0e6a61642053d4ed9043b6b0e1124}.
\bibitem{mv} H. L. Montgomery and R. C. Vaughan, \emph{Multiplicative Number Theory I: Classical Theory}, Cambridge University Press, 2007, Chapter 13 (conditional estimates).
\bibitem{products} \emph{On the existence of products of primes in arithmetic progressions}, Bulletin of the London Mathematical Society, DOI \href{https://doi.org/10.1112/blms.12990}{10.1112/blms.12990}. Background on the distinct problem of small-prime products in reduced residue classes.
\bibitem{zerosum} S. J. Lee and J. S. Oh, On zero-sum free sequences contained in random subsets of finite cyclic groups, \emph{Discrete Applied Mathematics} 328 (2023), 118--127. DOI \href{https://doi.org/10.1016/j.dam.2023.01.009}{10.1016/j.dam.2023.01.009}. Background on zero-sum-free sequences, not the allocation theorem here.
\end{thebibliography}
\end{document}
